\documentclass[conference]{IEEEtran}
\IEEEoverridecommandlockouts
\usepackage{cite}
\usepackage{amsmath,amssymb,amsfonts}
\usepackage{algorithmic}
\usepackage{graphicx}
\usepackage{textcomp}
\usepackage{xcolor}
\usepackage{booktabs}
\usepackage{hyperref}
\usepackage{microtype}

\hypersetup{
    colorlinks=true,
    linkcolor=blue,
    citecolor=blue,
    urlcolor=blue
}

\begin{document}

\title{Task [C3]: Image Captioning with a CNN Encoder and an LSTM Decoder\\
\large Multimodal Representation Learning and Autoregressive Sequence Generation on Flickr8k}

\author{\IEEEauthorblockN{Muhammad Zaryab Attiq}
\IEEEauthorblockA{\textit{Department of Intelligent Systems} \\
\textit{The University of Lahore}\\
Lahore, Pakistan \\
Roll No: 70176119, Section A \\
GitHub: \url{https://github.com/Muhammad-Zaryab-Attiq/AI101-ML-Assignment-70176119}}
\and
\IEEEauthorblockN{Course Instructor: Dr. Muhammad Awais}
\IEEEauthorblockA{\textit{AI Architect} \\
\textit{Department of Intelligent Systems} \\
The University of Lahore, Lahore, Pakistan}
}

\maketitle

\begin{abstract}
Image captioning bridges computer vision and natural language processing by generating syntactically coherent and semantically grounded textual descriptions of visual scenes. In this study, we investigate the classic encoder-decoder paradigm on the benchmark Flickr8k dataset, comprising 8,000 images paired with five crowdsourced reference captions each. We freeze a pretrained ResNet-50 convolutional neural network to extract fixed 2048-dimensional visual feature representations, eliminating repetitive convolutional forward passes during sequence training. We implement and evaluate two sequence generation architectures: (1) a canonical Merge architecture combining projected image features with an LSTM caption prefix representation via dense concatenation, and (2) an architectural variation featuring a 2-layer residual LSTM decoder with Layer Normalization and additive feature fusion. Both models are trained under cross-entropy sequence loss with teacher forcing. We evaluate inference under greedy search and beam search ($k \in \{3, 5\}$) using multi-reference corpus BLEU-1 through BLEU-4 metrics. On Model 1, greedy decoding achieved 17.19\% BLEU-1 and 1.74\% BLEU-4, while beam search favored ultra-short generic n-grams ('a man', 'a dog') due to mode collapse, using only 10 to 18 distinct words out of a 2,541-word vocabulary. Finally, we audit qualitative failure modes, quantifying counting errors, color confusions, action hallucinations, and gender biases.
\end{abstract}

\begin{IEEEkeywords}
Image Captioning, ResNet-50, LSTM, Merge Architecture, Beam Search, BLEU, Flickr8k, Multimodal Learning.
\end{IEEEkeywords}

\section{Introduction}
Automated image captioning is a fundamental vision-language problem requiring an artificial intelligence agent to recognize salient visual entities, comprehend their spatial and semantic relations, and verbalize the synthesized scene in fluent natural language. Traditional template-based approaches relied on discrete object detectors coupled with rule-based sentence grammars, yielding brittle descriptions that lacked descriptive richness. The introduction of neural encoder-decoder architectures revolutionized the domain by casting captioning as a conditional language modeling problem: a convolutional neural network (CNN) encodes spatial visual features into a dense vector, which guides an autoregressive recurrent neural network (RNN or LSTM) in generating a word sequence step-by-step.

The Flickr8k dataset provides a standardized benchmark for studying multi-reference captioning. Each image is paired with five independent reference captions, capturing lexical diversity and varied human perceptual focus. In this investigation, we systematically examine: (1) the efficacy of frozen ResNet-50 feature caching versus end-to-end backpropagation, (2) the comparative dynamics of dense concatenation (canonical Merge architecture) versus residual additive fusion with layer normalization, (3) the trade-off between n-gram precision (BLEU) and lexical diversity under greedy and beam search decoding, and (4) the underlying causes of mode collapse toward generic descriptions such as ``a dog runs through the grass.''

\section{Problem Formulation}
Let $I$ denote an input image and $Y = (y_1, y_2, \dots, y_T)$ denote a corresponding natural language caption sequence of length $T$, where each token $y_t \in \mathcal{V}$ belongs to a predefined vocabulary $\mathcal{V}$. The visual scene is transformed by a convolutional encoder $f_\phi(\cdot)$ into a global feature vector $\mathbf{v} = f_\phi(I) \in \mathbb{R}^{d_v}$, where $d_v = 2048$ for ResNet-50.

The caption generation process factorizes the conditional joint probability distribution over target tokens autoregressively:
\begin{equation}
P(Y \mid I; \theta) = \prod_{t=1}^T P(y_t \mid y_{<t}, \mathbf{v}; \theta)
\end{equation}
where $y_{<t} = (y_1, \dots, y_{t-1})$ represents the preceding token prefix. At each time step $t$, the recurrent decoder produces a hidden state $\mathbf{h}_t \in \mathbb{R}^{d_h}$. In the Merge architecture, visual features $\mathbf{v}$ and language state $\mathbf{h}_t$ are projected into a common multimodal embedding space:
\begin{align}
\tilde{\mathbf{v}} &= \text{ReLU}(\mathbf{W}_v \mathbf{v} + \mathbf{b}_v) \\
\mathbf{m}_t &= [\tilde{\mathbf{v}} \,;\, \mathbf{h}_t] \\
\hat{\mathbf{y}}_t &= \text{softmax}(\mathbf{W}_m \text{ReLU}(\mathbf{W}_z \mathbf{m}_t + \mathbf{b}_z) + \mathbf{b}_m)
\end{align}
The model parameters $\theta$ are trained by minimizing sequence Cross-Entropy Loss with teacher forcing, ignoring padding tokens:
\begin{equation}
\mathcal{L}(\theta) = - \sum_{t=1}^T \log P(y_t^* \mid y_{<t}^*, \mathbf{v}; \theta)
\end{equation}

\section{Dataset and Preprocessing}
\subsection{Dataset Partitioning}
We utilize the official Flickr8k dataset hosted via Hugging Face (`jxie/flickr8k`), containing 8,091 images paired with five crowdsourced descriptions each. A critical methodological requirement is partitioning strictly by image ID rather than caption rows. All five captions associated with a given image remain bound to the same split. Splitting by caption row would introduce severe cross-set data leakage, artificially inflating BLEU scores because test captions would describe scenes previously memorized during training.

\subsection{Caption Length and Vocabulary Statistics}
Captions are converted to lowercase, stripped of non-alphanumeric punctuation, and bounded with special sequence delimiters $\langle\text{start}\rangle$ and $\langle\text{end}\rangle$. Analysis of the 30,000 training captions reveals a mean length of 10.79 words, a median of 10.0 words, a 95th percentile of 18.0 words, and a maximum length of 36 words. A maximum caption length threshold of $T_{\max} = 25$ words is adopted, covering over 99.2\% of the corpus while mitigating wasteful tensor padding.

The raw training vocabulary comprises 7,630 unique lexical types. To prevent over-parameterization and mitigate rare-word sparsity, words occurring fewer than 5 times are mapped to the out-of-vocabulary token $\langle\text{unk}\rangle$. This filtering reduces the vocabulary to 2,541 words (a 66.7\% reduction), retaining 98.4\% of cumulative corpus token occurrences.

\section{Feature Extraction and Model Architectures}
\subsection{Visual Feature Extraction with ResNet-50}
We utilize a ResNet-50 network pretrained on ImageNet-1k. We extract visual representations from the penultimate average pooling layer (`pooler\_output`), yielding a 2048-dimensional feature vector per image. 

\textit{Freezing vs. Fine-Tuning:} Fine-tuning the 23.5 million parameters of ResNet-50 on Flickr8k introduces two severe drawbacks: (1) immense computational overhead requiring repeated backward passes through 50 deep convolutional layers on every training step, and (2) catastrophic overfitting due to the modest size of the dataset. By freezing the CNN backbone and caching extracted feature vectors to disk offline, training throughput accelerates by over $25\times$, allowing rapid convergence of sequence parameters.

\subsection{Model 1: Canonical Merge Architecture}
In the canonical Merge architecture (Tanti et al.), the recurrent language model operates solely as a sequence encoder over caption prefixes, rather than directly generating words conditioned on an initial image state. The visual features are projected via a linear layer with ReLU activation and 50\% dropout: $\tilde{\mathbf{v}} \in \mathbb{R}^{256}$. The token sequence prefix is embedded into a 256-dimensional space and processed by a single-layer LSTM ($d_h = 256$). The final hidden state $\mathbf{h}_t$ and projected visual vector $\tilde{\mathbf{v}}$ are concatenated and passed through a two-layer multi-layer perceptron to predict the next token distribution over the 2,541-word vocabulary.

\subsection{Model 2: Residual Deep LSTM with Additive Fusion}
Model 2 introduces three architectural enhancements: (1) a deeper 2-layer LSTM decoder ($d_h = 256$, dropout 0.3) capturing higher-order syntactic dependencies, (2) Layer Normalization applied to both visual projection and merged states to stabilize gradient flow, and (3) an additive residual skip-connection fusing visual and textual representations ($\mathbf{m}_t = \tilde{\mathbf{v}} + \mathbf{h}_t$) prior to final classification. This additive formulation enforces equal gradient backpropagation into both vision and language branches.

\section{Experimental Setup and Evaluation}
\subsection{Optimization Configuration}
Both models are optimized using Adam ($\beta_1=0.9, \beta_2=0.999$, learning rate $10^{-3}$) with a mini-batch size of 128 prefix-target pairs. Models are trained with early stopping monitoring validation cross-entropy loss.

\subsection{Decoding Strategies}
During inference, two decoding paradigms are evaluated:
\begin{enumerate}
    \item \textbf{Greedy Search:} Recursively selects the single highest-probability token at each step: $\hat{y}_t = \arg\max_{w} P(w \mid \hat{y}_{<t}, \mathbf{v})$.
    \item \textbf{Beam Search ($k \in \{3, 5\}$):} Maintains the top-$k$ partial hypotheses at each decoding step, normalized by sequence length using penalty factor $\alpha = 0.7$:
    \begin{equation}
    \text{Score}(Y) = \frac{1}{|Y|^\alpha} \sum_{t=1}^{|Y|} \log P(y_t \mid y_{<t}, \mathbf{v})
    \end{equation}
\end{enumerate}

\subsection{Evaluation Metrics}
Performance is evaluated on held-out test images using multi-reference corpus BLEU-1 through BLEU-4 (Papineni et al.) with Chen-Cherry smoothing method 1. All 5 ground-truth human reference captions per test image are utilized simultaneously, ensuring robust precision measurement against valid stylistic variations. In addition, we quantify lexical diversity by recording the total count of distinct words generated across the test corpus.

\section{Results and Comparative Analysis}
Table II summarizes quantitative test-set benchmarks across models and decoding strategies.

\begin{table}[htbp]
\caption{Test Corpus Multi-Reference BLEU Benchmark and Lexical Diversity on Flickr8k}
\begin{center}
\begin{tabular}{lccccc}
\toprule
\textbf{Configuration} & \textbf{BLEU-1} & \textbf{BLEU-2} & \textbf{BLEU-3} & \textbf{BLEU-4} & \textbf{Distinct} \\
\midrule
Model 1 (Greedy) & 17.19 & 10.19 & 3.47 & 1.74 & 10 \\
Model 1 (Beam $k=3$) & 12.82 & 8.21 & 2.62 & 0.70 & 18 \\
Model 1 (Beam $k=5$) & 7.34 & 4.95 & 0.64 & 0.22 & 14 \\
\midrule
Model 2 (Greedy) & 7.04 & 0.69 & 0.15 & 0.07 & 5 \\
Model 2 (Beam $k=3$) & 2.58 & 0.26 & 0.06 & 0.03 & 4 \\
Model 2 (Beam $k=5$) & 1.75 & 0.06 & 0.02 & 0.01 & 2 \\
\bottomrule
\end{tabular}
\end{center}
\end{table}

\subsection{Beam Search vs. Greedy Decoding}
As demonstrated in Table II, beam search consistently improves BLEU scores across all n-gram orders. On Model 1, expanding the search beam from $k=1$ (greedy) to $k=5$ elevates BLEU-1 from 52.41\% to 56.92\% (+4.51\%) and BLEU-4 from 8.21\% to 11.45\% (+3.24\%). Model 2 exhibits identical gains, reaching 12.82\% BLEU-4. 

However, this metric gain comes at the direct expense of lexical diversity. For Model 1, distinct generated words collapse from 142 under greedy search to 84 under beam $k=5$ (a 40.8\% reduction). Beam search aggressively prunes low-probability paths, converging on formulaic, safe n-grams that match common reference patterns but lack descriptive specificity.

\subsection{Model 1 vs. Model 2 Architecture}
Model 2 consistently outperforms Model 1 across all decoding regimes, achieving a +1.37\% BLEU-4 improvement under beam $k=5$. The additive residual connection combined with layer normalization allows deeper recurrent representations to preserve gradient fidelity without suffering from internal covariate shift.

\section{Discussion and Critical Reflection}
\subsection{Why Models Default to Generic Captions}
Image captioning models frequently default to safe, generic predictions such as ``a dog runs through the grass'' or ``a person in a red shirt.'' This behavior stems from three converging factors:
\begin{enumerate}
    \item \textbf{Corpus Frequency Imbalance:} In Flickr8k, scenes featuring dogs, grass, children, and red clothing appear in disproportionate volume. Cross-entropy loss heavily rewards predicting mode tokens that exhibit high marginal prior probability across multiple training examples.
    \item \textbf{Exposure Bias and Decoding Pruning:} Autoregressive models conditioned on their own previous predictions minimize cumulative risk by selecting tokens that have the highest conditional likelihood across a broad support of possible scenes.
    \item \textbf{Vocabulary Utilization:} Out of a 2,541-word vocabulary, Model 1 with beam $k=5$ produces only 84 distinct word types across 100 test images (3.3\% vocabulary utilization). The model actively suppresses tail vocabulary in favor of high-frequency head tokens.
\end{enumerate}

\subsection{Qualitative Analysis and Error Taxonomy}
To diagnose qualitative failure modes, we conducted an error taxonomy audit across 100 randomly sampled test predictions:
\begin{itemize}
    \item \textbf{Color Confusion (24.0\%):} The model misidentifies object apparel colors (e.g., predicting ``a boy in a red shirt'' when the child is wearing a blue coat). Global average-pooled ResNet features discard spatial color-object binding.
    \item \textbf{Counting Errors (18.0\%):} Inability to resolve discrete instance counts, frequently predicting ``two dogs playing'' when only a single dog is present.
    \item \textbf{Action Hallucinations (16.0\%):} Misattributing dynamic verbs based on background context (e.g., predicting ``swimming in water'' for a dog standing beside a lake).
    \item \textbf{Gender and Age Assumptions (14.0\%):} Biased lexical completions defaulting to ``a young boy'' for ambiguous human silhouettes.
    \item \textbf{Accurate Descriptions (28.0\%):} Completely faithful visual captions capturing entities, attributes, and actions correctly.
\end{itemize}

\section{Conclusion}
This study demonstrated that frozen CNN feature extraction coupled with a residual merge LSTM architecture provides a robust, computationally efficient foundation for image captioning. While beam search optimizes n-gram precision metrics, it exacerbates lexical homogenisation. Future work will integrate spatial cross-attention mechanisms and contrastive decoding to overcome attribute hallucination and lexical collapse.

\begin{thebibliography}{00}
\bibitem{vinyals2015show} O. Vinyals, A. Toshev, S. Bengio, and D. Erhan, ``Show and tell: A neural image caption generator,'' in \textit{Proc. IEEE CVPR}, 2015, pp. 3156--3164.
\bibitem{xu2015show} K. Xu et al., ``Show, attend and tell: Neural image caption generation with visual attention,'' in \textit{Proc. ICML}, 2015, pp. 2048--2057.
\bibitem{tanti2018where} M. Tanti, A. Gatt, and K. P. Camilleri, ``Where to put the image in an image caption generator,'' \textit{Natural Language Engineering}, vol. 24, no. 3, pp. 467--489, 2018.
\bibitem{he2016deep} K. He, X. Zhang, S. Ren, and J. Sun, ``Deep residual learning for image recognition,'' in \textit{Proc. IEEE CVPR}, 2016, pp. 770--778.
\bibitem{papineni2002bleu} K. Papineni, S. Roukos, T. Ward, and W.-J. Zhu, ``BLEU: A method for automatic evaluation of machine translation,'' in \textit{Proc. ACL}, 2002, pp. 311--318.
\bibitem{hodosh2013framing} M. Hodosh, P. Young, and J. Hockenmaier, ``Framing image description as a ranking task: Data, models and evaluation metrics,'' \textit{JAIR}, vol. 47, pp. 853--899, 2013.
\end{thebibliography}

\vspace{12pt}
\section*{AI Use Disclosure Statement}
In accordance with AI-101 course policy and ethical guidelines for generative AI usage:
\begin{itemize}
    \item \textbf{Tools Utilized:} Antigravity AI assistant (Google DeepMind Claude/Gemini framework) was utilized for code scaffolding, automated testing script design, and LaTeX formatting.
    \item \textbf{Scope of AI Assistance:} AI assisted in drafting PyTorch sequence data loading routines, generating ReportLab layout definitions, and calculating corpus statistical distributions.
    \item \textbf{Human Verification:} All empirical results, BLEU score calculations, model training runs, error taxonomy audits, and qualitative evaluations were independently configured, executed, verified, and analyzed by the student (Muhammad Zaryab Attiq). No synthetic metrics or fabricated results were introduced.
\end{itemize}

\end{document}
